\documentclass[11pt,a4paper]{ctexart}
\usepackage[margin=2.35cm]{geometry}
\usepackage{amsmath,amssymb,amsthm,mathtools,bm,mathrsfs}
\usepackage{xcolor,booktabs,longtable,array,enumitem,hyperref}
\hypersetup{colorlinks=true,linkcolor=blue!55!black,urlcolor=blue!55!black,citecolor=blue!55!black}
\setlist{nosep}
\setlength{\emergencystretch}{3em}

\newtheorem{theorem}{定理}[section]
\newtheorem{lemma}[theorem]{引理}
\newtheorem{corollary}[theorem]{推论}
\newtheorem{proposition}[theorem]{命题}
\theoremstyle{definition}
\newtheorem{definition}[theorem]{定义}
\newtheorem{remark}[theorem]{适用范围}
\newtheorem{warning}[theorem]{警告}

\newcommand{\F}{\mathbb F}
\newcommand{\C}{\mathbb C}
\newcommand{\T}{\mathbb T}
\newcommand{\G}{\F_p^\times}
\newcommand{\Gh}{\widehat{\G}}
\newcommand{\ee}{\mathrm e}
\newcommand{\ep}{e_p}
\newcommand{\1}{\mathbf 1}
\newcommand{\Tr}{\operatorname{Tr}}
\newcommand{\diag}{\operatorname{diag}}
\newcommand{\Kl}{\operatorname{Kl}}
\newcommand{\supp}{\operatorname{supp}}
\newcommand{\eps}{\varepsilon}
\newcommand{\Res}{\operatorname*{Res}}

\title{GFT 可复用工具箱与优美恒等式\\
\large 有限域、有限群与算子层的严格版本及适用范围}
\author{从十五份研究日志中筛选并重新证明}
\date{2026年8月28日}

\begin{document}
\maketitle
\sloppy

\begin{abstract}
本文只保留能够在给定定义下完整证明的 GFT 性质，不保留任何依赖未证
Goldbach 接口的“候选定理”。核心包括：有限 hit-count 的精确 DFT projector；
unitary Mellin 的方向审计；sum--product joint-charge transform 的 exact Gram；
reciprocal fibre 与 projective map；Cayley 变量把有理 complete sum精确化为
Kloosterman sum；Gauss--Jacobi coboundary 与循环 telescoping；Gauss-square
序列的 multiplicative Fourier/Kloosterman 对偶；centered Parseval；
fibrewise dilation 的酉性与 Schatten ideal property；以及“一般 entry shear
不保持 Schatten 范数”的反例。修订版还补入 smooth HWAC 的逐 Fourier-mode
Cauchy residue 与 High-Conductor ZOT 的完整 Euler-product/rational-sampling
证明，并明确删除 sharp-cutoff 与整体全纯闭合的错误版本。本文进一步补入
gcd-carried FSAC、origin-labelled tensor reconstruction、interval autocorrelation、
cross-prime character incidence、BP trace 的 sum--product pushforward 与 divisor
rigidity。本版还重新证明 Exact Parity--Gauge Decomposition、finite ModelGlue、
factor-scale power-sum reconstruction、有限 alternating endpoint jet，以及
RSC/Bellman 的 finite-witness 与 supersolution comparison。每条定理都标明 characteristic、非主角色、非零分母、measure、
normalization 与外部输入。
\end{abstract}

\section{统一约定}

除非另有说明，$p$ 为奇素数，$\F_p$ 为 $p$ 元域，
$\G=\F_p^\times$，$Q=|\G|=p-1$，$\Gh$ 为所有乘法角色。
取非平凡加法角色
\[
\ep(x)=\exp(2\pi i x/p).
\]
所有有限集合上的 $\ell^2$ 默认使用 counting measure。

加法 unitary Fourier 变换定义为
\[
\widehat f(h)=p^{-1/2}\sum_{x\in\F_p}f(x)\ep(-hx),
\tag{1.1}
\]
乘法 unitary Mellin 变换定义为
\[
(\mathcal M f)(\chi)=Q^{-1/2}\sum_{x\in\G}f(x)\overline{\chi(x)}.
\tag{1.2}
\]
其逆为
\[
f(x)=Q^{-1/2}\sum_{\chi\in\Gh}(\mathcal Mf)(\chi)\chi(x).
\tag{1.3}
\]

基础正交关系为
\begin{align}
\sum_{h\in\F_p}\ep(hx)&=p\,\1_{x=0},                 \tag{1.4}\\
\sum_{\chi\in\Gh}\chi(x)&=Q\,\1_{x=1}\qquad(x\in\G). \tag{1.5}
\end{align}

\section*{全局附录 A：HWAC 可严格使用的留数位移变换}
\addcontentsline{toc}{section}{全局附录 A：HWAC 可严格使用的留数位移变换}

本节不需要有限域。令 $\T=\mathbb R/\mathbb Z$，仍写
$e(t)=e^{2\pi it}$。设 $\mathfrak M^-\subset\mathfrak M^+\subset\T$，
$\mathfrak M^-$ 至多为 $J$ 个区间之并，且二者边界间距至少 $3\eta$。

\begin{theorem}[smooth Hardy--Wiener Arc Closure]
存在 $C^\infty$ partition
\(
\chi_{\mathfrak M}+\psi_{\mathfrak m}=1
\)，其中 $\chi=1$ 于 $\mathfrak M^-$、
$\supp\chi\subset\mathfrak M^+$，并且
\[
\psi_{\mathfrak m}(\alpha)=\sum_{r\in\mathbb Z}\omega_re(r\alpha),
\qquad
\sum_r|\omega_r|\ll1+J\left(1+\log\frac2\eta\right).   \tag{A.1}
\]
若 $F(\alpha)=\sum_xa_xe(x\alpha)$、
$G(\alpha)=\sum_yb_ye(y\alpha)$ 为有限指数和，则
\[
\boxed{
\int_\T\psi_{\mathfrak m}(\alpha)F(\alpha)G(\alpha)e(-N\alpha)d\alpha
=\sum_r\omega_r\sum_{x+y=N-r}a_xb_y.}                 \tag{A.2}
\]
而且每个 shifted coefficient 都有严格 residue 表示
\[
\sum_{x+y=N-r}a_xb_y
=\Res_{w=0}\frac{F(w)G(w)}{w^{N-r+1}}.                 \tag{A.3}
\]
\end{theorem}

\begin{proof}
取 $\rho\ge0$、$\rho\in C_c^\infty((-1,1))$、$\int\rho=1$。把
$\mathfrak M^-$ 扩张 $\eta$ 为 $E$，令
$\chi=\1_E*\rho_\eta$、$\psi=1-\chi$。间距假设给所需 support 性质。
对 $r\ne0$，一个区间的 Fourier coefficient 绝对值至多
$1/(\pi|r|)$；而两次分部积分给
$|\widehat\rho(\eta r)|\ll(1+\eta|r|)^{-2}$。故
\[
\sum_{r\ne0}|\widehat\chi(r)|
\ll J\sum_{r\ge1}\frac1r(1+\eta r)^{-2}
\ll J\left(1+\log\frac2\eta\right),
\]
证明 (A.1) 与 Fourier 级数的一致绝对收敛。于是可交换有限的 $x,y$-和、
绝对可和的 $r$-和与积分；$\int_\T e(k\alpha)d\alpha=\1_{k=0}$ 立即给
(A.2)。最后 (A.3) 是 Cauchy coefficient formula，且 residue 与 $r$-和的
交换仍由 (A.1) 保证。
\end{proof}

\begin{warning}[适用范围]
留数是对 (A.2) 的\emph{每个 Fourier mode} 作 coefficient extraction。
一般 $C^\infty$ cutoff 不必在圆环中全纯；所以把整个正/负频率部分分别
向内/向外闭合不是本定理。sharp interval indicator 的系数为 $O(1/|r|)$，
Wiener $\ell^1$ 发散，故 sharp-HWAC 版本不能调用。
\end{warning}

\section*{全局附录 B：ZOT 把 zero orbit 转移到奇异级数尾部}
\addcontentsline{toc}{section}{全局附录 B：ZOT 把 zero orbit 转移到奇异级数尾部}

固定 $W_1,W_2\in C_c^\infty((0,1))$，定义
\[
\mathfrak J_W(\lambda)=\int W_1(t)W_2(\lambda-t)dt,
\quad
\mathfrak S(M)=\sum_{q\ge1}\frac{\mu(q)^2}{\varphi(q)^2}c_q(-M). \tag{B.1}
\]
使用 $\Lambda=\mu*\log$ 后，对 $du+ev=M$ 的解参数作一维 Poisson。
zero frequency 的格点密度为 $(d,e)/(de)$，故其算术核是
\[
D_M(X,Y)=\sum_{\substack{d,e\ge1\\(d,e)\mid M}}
\mu(d)\mu(e)\frac{(d,e)}{de}
\log_+(X/d)\log_+(Y/e).                                \tag{B.2}
\]

\begin{theorem}[Möbius zero-orbit resummation]
若 $M,X,Y\asymp N$，则对任意 $A>0$，
\[
\boxed{D_M(X,Y)=\mathfrak S(M)+O_A((\log N)^{-A}).}     \tag{B.3}
\]
从而完整 Poisson zero orbit 为
\[
\boxed{Z_0(M)=N\mathfrak S(M)\mathfrak J_W(M/N)
+O_A(N(\log N)^{-A}).}                                 \tag{B.4}
\]
\end{theorem}

\begin{proof}
Perron 公式 $\log_+x=(2\pi i)^{-1}\int_{(c)}x^sds/s^2$ 给
\[
D_M(X,Y)=\frac1{(2\pi i)^2}\iint
\mathcal A_M(s,t)X^sY^t\frac{ds\,dt}{s^2t^2},          \tag{B.5}
\]
其中
\[
\mathcal A_M(s,t)=\sum_{\substack{d,e\ge1\\(d,e)\mid M}}
\mu(d)\mu(e)\frac{(d,e)}{d^{1+s}e^{1+t}}.              \tag{B.6}
\]
对 $p\nmid M$，Möbius local states 给
$1-p^{-1-s}-p^{-1-t}$；对 $p\mid M$，双重状态被允许并给
$1-p^{-1-s}-p^{-1-t}+p^{-1-s-t}$。因此
\[
\mathcal A_M(s,t)=\frac{\mathcal H_M(s,t)}
{\zeta(1+s)\zeta(1+t)},                                \tag{B.7}
\]
\begin{align*}
\mathcal H_M(s,t)
={}&\prod_{p\nmid M}\frac{1-p^{-1-s}-p^{-1-t}}
{(1-p^{-1-s})(1-p^{-1-t})}\\
&\times\prod_{p\mid M}
\frac{1-p^{-1-s}-p^{-1-t}+p^{-1-s-t}}
{(1-p^{-1-s})(1-p^{-1-t})}.                            \tag{B.8}
\end{align*}
第一乘积局部为 $1+O(p^{-2+o(1)})$，故在原点邻域解析。若 $M$ 偶，
$p=2$ 给 $2$；$p>2,p\nmid M$ 给 $1-(p-1)^{-2}$；$p\mid M$ 相对基线
给 $(p-1)/(p-2)$。若 $M$ 奇，$p=2$ 因子为 $0$。所以恰有
\[
\mathcal H_M(0,0)=\mathfrak S(M).                      \tag{B.9}
\]
取 $T=e^{\sqrt{\log N}}$，把 (B.5) 截断在 $|\Im s|,|\Im t|\le T$，
并从 $1/\log N$ 移到 $-\kappa/\log T$。经典 zeta 零点自由区保证矩形内
无 zeta 零点，$1/\zeta(1+s)\ll\log^{O(1)}T$，而
$\mathcal H_M\ll\log^{O(1)}N$。由
$1/\zeta(1+s)=s+O(s^2)$，穿过的双 residue 是 (B.9)。新直线贡献
$\ll\log^{O(1)}N e^{-c\sqrt{\log N}}$；水平边和原直线尾部为
$\ll T^{-1}\log^{O(1)}N$。这证明 (B.3)。在 zero-frequency integral 中，
$X=x,Y=M-x\asymp N$；代入 (B.3) 并令 $x=Nt$ 即得 (B.4)。
\end{proof}

现在取 $P=(\log N)^B$，major core/envelope 的宽度分别为
$R/(qN)$、$2R/(qN)$，其中 $R=P(\log N)^C$；令
$\psi_{\mathfrak m}=\sum_r\omega_re(r\alpha)$ 为对应 HWAC cutoff。

\begin{lemma}[rational sampling]
对 $J\in C_c^\infty(\mathbb R)$，置
\[
T_m=\sum_r\omega_rJ(1-r/N)c_m(-(N-r)).                 \tag{B.10}
\]
任意 $L>0$ 下，
\[
T_m\ll\varphi(m)(1+R/m)^{-L}\quad(m\le P),             \tag{B.11}
\]
而当 $P<m\le N/(4R)$ 时，
\[
T_m=J(1)c_m(-N)
+O(\varphi(m)(1+R/P)^{-L}).                            \tag{B.12}
\]
\end{lemma}

\begin{proof}
展开 $c_m$。Fourier inversion 给
\[
\sum_r\omega_rJ(1-r/N)e(ar/m)
=\int\widehat J(t)e(t)\psi_{\mathfrak m}(a/m-t/N)dt.  \tag{B.13}
\]
若 $m\le P$，$a/m$ 在 major core，离 $\psi$ 的 support 至少
$\gg R/(mN)$；利用 $\widehat J$ 的任意次幂衰减得到 (B.11)。若
$P<m\le N/(4R)$，任何 reduced $b/q,q\le P$ 都满足
$|a/m-b/q|\ge1/(mq)\ge4R/(qN)$，所以 $a/m$ 在 minor plateau 且到
envelope 距离 $\gg R/(PN)$。在 (B.13) 中以 $1$ 替换 $\psi$；主项为
$J(1)$，尾积分给 (B.12)。
\end{proof}

\begin{theorem}[High-Conductor ZOT]
令
\[
\mathfrak S_P(N)=\sum_{m\le P}\frac{\mu(m)^2}{\varphi(m)^2}c_m(-N).
\]
则 minor zero orbit 满足
\[
\boxed{Z_{\mathfrak m}^{(0)}(N)
=N\mathfrak J_W(1)(\mathfrak S(N)-\mathfrak S_P(N))
+O_A(N(\log N)^{-A}).}                                 \tag{B.14}
\]
\end{theorem}

\begin{proof}
由 (A.2)、(B.4)，左侧除 $O_A(N\log^{-A}N)$ 外等于
$N\sum_r\omega_r\mathfrak S(N-r)\mathfrak J_W(1-r/N)$。
在 (B.1) 的 $m$-和中分成 $m\le P$、$P<m\le Y=N/(4R)$、$m>Y$。
严格地说，先截断在 $m\le Q$；下述界对 $Q$ 一致，再用 (B.15)
令 $Q\to\infty$，所以没有交换条件收敛级数的问题。
前两段分别用 (B.11)、(B.12)。误差求和只需
$\sum_{m\le X}\mu(m)^2/\varphi(m)\ll\log X$；它由
$n/\varphi(n)=\sum_{d\mid n}\mu(d)^2/\varphi(d)$ 直接得到。
最后一段调用 Montgomery--Vaughan 点态界
\[
|\widetilde{\mathfrak S}_Y(k)|
\ll d(k)(\log\log3Y)^2/Y.                              \tag{B.15}
\]
在 Archimedean support 上 $k=N-r\asymp N$；结合
$d(k)=N^{o(1)}$、$Y=N/\log^{O(1)}N$ 与 (A.1)，该段为
$O_A(\log^{-A}N)$。于是中段可延伸到 $m=\infty$，得到 (B.14)。
\end{proof}

\begin{remark}
ZOT 的结论不是 zero orbit 很小，而是它精确承载 truncated major arcs 缺少的
high-conductor singular-series tail。若主弧外部输入为
$N\mathfrak J_W(1)\mathfrak S_P(N)+O_A(N\log^{-A}N)$，两者相加才得到
$N\mathfrak J_W(1)\mathfrak S(N)+o(N)$。
\end{remark}

\section{有限 hit-count 的精确 DFT projector}

\begin{theorem}[有限计数选择器]
设 $K\ge1$，$r,j\in\{0,1,\ldots,K-1\}$。则
\[
\boxed{
\1_{r=j}=\frac1K\sum_{\ell=0}^{K-1}
\exp\!\left(\frac{2\pi i\ell(r-j)}K\right).}
\tag{2.1}
\]
\end{theorem}

\begin{proof}
若 $r=j$，每项为 $1$。若 $r\ne j$，右侧是公比不为 $1$ 而 $K$ 次方为
$1$ 的几何级数，故和为 $0$。
\end{proof}

\begin{remark}
若一个算术计数 $R(n)$ 已先验满足 $0\le R(n)<K$，则 (2.1) 精确选择
$R(n)=j$。若只知道 $R(n)$ 是整数而无此范围，右侧选择的是
$R(n)\equiv j\pmod K$；这两个陈述不能混用。
\end{remark}

\section{Mellin 方向与 reciprocal fibre}

\begin{lemma}[unitary Mellin 的逆向]
若
\[
W(y)=Q^{-1/2}\sum_{D\in\Gh}\overline{D(y)}\widetilde W(D),
\]
则
\[
\boxed{\widetilde W(D)=Q^{-1/2}\sum_{y\in\G}D(y)W(y).}
\tag{3.1}
\]
\end{lemma}

\begin{proof}
代入右侧并交换求和，内层
$Q^{-1}\sum_y D(y)\overline{D'(y)}=\1_{D=D'}$。
\end{proof}

\begin{theorem}[Mellin--Jacobi reciprocal support]
取 $z,z',y,y'\in\G$ 且 $z\ne z'$，令
\[
r=z/z',\qquad a=(r-1)/r,\qquad b=1-r.
\]
若 $x\in\F_p\setminus\{0,1\}$ 满足
\[
yax=1,\qquad y'b(1-x)=1,
\tag{3.2}
\]
则
\[
\boxed{z(y^{-1}-1)=z'(y'^{-1}-1).}
\tag{3.3}
\]
\end{theorem}

\begin{proof}
因 $a=(z-z')/z$、$b=(z'-z)/z'$，由 (3.2)
\[
x=\frac{z}{y(z-z')},\qquad
1-x=-\frac{z'}{y'(z-z')}.
\]
相加得到 $z/y-z'/y'=z-z'$，整理即 (3.3)。
\end{proof}

\begin{corollary}[projective charge map]
令 $q=z(y^{-1}-1)$、$w=yz$。则 $z+q\ne0$，且
\[
\boxed{y=\frac{z}{z+q},\qquad
w=\frac{z^2}{z+q}.}
\tag{3.4}
\]
反过来，对 $z\in\G$、$q\ne-z$，(3.4) 唯一确定 $y,w\in\G$。
\end{corollary}

\begin{proof}
$z+q=z/y\ne0$，其余为直接代数运算。反向代入可验证。
\end{proof}

\begin{remark}
这一定理说明，在 (3.1) 的 convention 下，后续几何是
$z\mapsto z^2/(z+q)$，不是 affine translation $z\mapsto z+q$。
若改用相反 Mellin convention，coefficient formula 也必须同步改写。
\end{remark}

\section{sum--product joint-charge transform}

令 $\mathscr H=\ell^2(\G\times\G)$。定义 swap 与 diagonal projection：
\[
(\mathsf SF)(x,z)=F(z,x),\qquad
(\mathsf DF)(x,z)=\1_{x=z}F(x,z).
\tag{4.1}
\]
先定义 exact sum--product charge transform
\[
(\mathcal AF)(R,P)
=\sum_{\substack{x,z\in\G\\x+z=R,\ xz=P}}F(x,z),
\quad R\in\F_p,\ P\in\G,
\tag{4.2}
\]
再定义它在 $P$ 坐标上的 unitary character transform
\[
(\mathcal JF)(R,\eta)
=Q^{-1/2}\sum_{\substack{x,z\in\G\\x+z=R}}
F(x,z)\eta(xz),
\quad R\in\F_p,\ \eta\in\Gh.
\tag{4.3}
\]
于是
\(
\mathcal JF(R,\eta)=Q^{-1/2}\sum_{P\in\G}(\mathcal AF)(R,P)\eta(P)
\)；这里使用的是与 (1.2) 相差角色反演 $\eta\mapsto\eta^{-1}$ 的
unitary convention。

\begin{theorem}[exact sum--product Gram]
在 counting measure 下，
\[
\boxed{\mathcal A^*\mathcal A
=\mathcal J^*\mathcal J=I+\mathsf S-\mathsf D.}
\tag{4.4}
\]
特别地
\[
\boxed{\|\mathcal A\|_{2\to2}
=\|\mathcal J\|_{2\to2}\le\sqrt2.}
\tag{4.5}
\]
\end{theorem}

\begin{proof}
对 $F,G\in\mathscr H$，展开内积：
\[
\langle\mathcal JF,\mathcal JG\rangle
=\frac1Q\sum_{R,\eta}
\sum_{\substack{x+z=R\\x'+z'=R}}
F(x,z)\overline{G(x',z')}
\eta\!\left(\frac{xz}{x'z'}\right).
\]
由 (1.5)，角色和恰好强迫 $xz=x'z'$。这也正是直接展开
$\langle\mathcal AF,\mathcal AG\rangle$ 所得的条件。同时有
$x+z=x'+z'$，故 $(x,z)$ 与 $(x',z')$ 是同一首一二次多项式
$T^2-RT+xz$ 的有序根；于是 $(x',z')=(x,z)$ 或 $(z,x)$。
当 $x=z$ 时两种选择重合，须减去一次 diagonal，得到 (4.4)。
由于 $\mathsf S$ 是酉自伴 involution、$\mathsf D$ 是正交投影，
$I+\mathsf S-\mathsf D\le2I$，故 (4.5)。
\end{proof}

\begin{corollary}[对称/反对称分解]
在对称子空间 $\mathscr H_+=\{F:\mathsf SF=F\}$ 上，
\[
I\le\mathcal J^*\mathcal J=2I-\mathsf D\le2I.
\tag{4.6}
\]
在反对称子空间上 $\mathcal J=0$。
\end{corollary}

\begin{proof}
对称情形代入 (4.4)。反对称函数在 $x=z$ 上为零，且同一无序根对的两项相消。
\end{proof}

\begin{remark}
此结论只给 $L^2$/Bessel assembly 与有界 fibre；它本身不产生
Schatten-$4$ 或 Schatten-$8$ 的 power saving。
\end{remark}

\section{projective Cayley 曲线就是 joint-charge 对角}

\begin{theorem}[Vi\`ete 对角字典]
映射
\[
r\in\F_p\setminus\{0,1\}
\longmapsto
(x,z)=\left(r^{-1},(1-r)^{-1}\right)
\tag{5.1}
\]
是到集合
\[
\Delta=\{(x,z)\in\G\times\G:x+z=xz\}
\]
的双射。并且交换 $(x,z)$ 对应 $r\mapsto1-r$。
\end{theorem}

\begin{proof}
直接计算
\[
x+z=\frac{1}{r}+\frac1{1-r}
=\frac1{r(1-r)}=xz.
\]
反之若 $x+z=xz$，令 $r=x^{-1}$，则
$1-r=1-x^{-1}=(x-1)/x=z^{-1}$，其中最后一步由
$z(x-1)=x$ 得到。故逆映射唯一。交换两坐标显然把 $r$ 变为 $1-r$。
\end{proof}

\begin{corollary}
在 $\mathcal A$ 的 exact-charge target
$\ell^2(\F_p\times\G)$ 上，限制 $R=P$ 是乘子
$\1_{R=P}$，因而是范数 $1$ 的正交投影。不能在
$\mathcal J$ 的角色坐标 $(R,\eta)$ 中直接写 $R=\eta$；若从
$\mathcal J$ 出发，必须先对 $\eta$ 作 unitary inverse Mellin。
\end{corollary}

\section{Cayley 变换与精确 Kloosterman 化}

定义经典 complete Kloosterman sum
\[
S(A,C;p)=\sum_{t\in\G}\ep(A/t+Ct).
\tag{6.1}
\]

\begin{theorem}[projective complete sum的 Kloosterman identity]
对任意 $A,C\in\F_p$，
\[
\boxed{
\sum_{r\in\F_p\setminus\{0,1\}}
\ep\!\left(\frac A r+\frac C{1-r}\right)
=\ep(A+C)S(A,C;p)-1.}
\tag{6.2}
\]
\end{theorem}

\begin{proof}
置 $t=r/(1-r)$。这是
$\F_p\setminus\{0,1\}$ 到 $\F_p\setminus\{0,-1\}$ 的双射，逆为
$r=t/(1+t)$。并且
\[
\frac1r=1+\frac1t,\qquad \frac1{1-r}=1+t.
\]
故左侧为
\[
\ep(A+C)\sum_{t\in\G\setminus\{-1\}}\ep(A/t+Ct).
\]
补上 $t=-1$ 的项，其相位为 $\ep(-A-C)$；乘回外因子后恰为 $1$，得到 (6.2)。
\end{proof}

\begin{corollary}[generic/degenerate 三分]
\[
\left|\sum_{r\ne0,1}\ep(A/r+C/(1-r))\right|
\le
\begin{cases}
2\sqrt p+1,&AC\ne0,\\
2,&\text{恰有一个 of }A,C\text{ 为 }0,\\
p-2,&A=C=0.
\end{cases}
\tag{6.3}
\]
\end{corollary}

\begin{proof}
$AC\ne0$ 时用 Weil 的 Kloosterman bound
$|S(A,C;p)|\le2\sqrt p$。若恰有一个为零，非平凡 additive character
在 $\G$ 上的和为 $-1$。若二者都为零，原和有 $p-2$ 项。
\end{proof}

\begin{remark}
(6.2) 是 exact identity；但把单个 entry 的 (6.3) 放入 Schatten trace 时，
仍需同一 tensor measure 上的 Plancherel/orthogonality。entrywise Weil
不能自动升级为 operator saving。
\end{remark}

\section{Gauss--Jacobi coboundary 与循环 telescoping}

对 $\chi\in\Gh$ 定义 Gauss sum
\[
\tau(\chi)=\sum_{x\in\G}\chi(x)\ep(x),
\tag{7.1}
\]
对乘法角色 $\alpha,\beta$ 定义 Jacobi sum
\[
J(\alpha,\beta)=\sum_{x\in\F_p}\alpha(x)\beta(1-x),
\tag{7.2}
\]
并按通常约定把角色在 $0$ 延拓为 $0$。

\begin{theorem}[normalized Jacobi 是 Gauss phase 的二余边界]
若 $\alpha,\beta,\alpha\beta$ 均非主角色，令
\[
\rho(\chi)=\frac{\tau(\chi)}{\sqrt p},
\qquad
\jmath(\alpha,\beta)=\frac{J(\alpha,\beta)}{\sqrt p}.
\]
则 $|\rho(\chi)|=1$，并且
\[
\boxed{
\jmath(\alpha,\beta)
=\rho(\alpha)\rho(\beta)\overline{\rho(\alpha\beta)}.}
\tag{7.3}
\]
\end{theorem}

\begin{proof}
经典 Gauss--Jacobi 恒等式给
\[
J(\alpha,\beta)=\frac{\tau(\alpha)\tau(\beta)}{\tau(\alpha\beta)}.
\]
非主角色满足 $|\tau(\chi)|=\sqrt p$，故
$1/\rho(\chi)=\overline{\rho(\chi)}$。除以 $\sqrt p$ 即得 (7.3)。
\end{proof}

\begin{corollary}[cocycle identity]
若下式涉及的每个角色均落在定理的非主适用域，则
\[
\boxed{
\jmath(\alpha,\beta)\jmath(\alpha\beta,\gamma)
=\jmath(\beta,\gamma)\jmath(\alpha,\beta\gamma).}
\tag{7.4}
\]
\end{corollary}

\begin{proof}
将 (7.3) 代入两边；所有中间 $\rho$ 因子望远镜相消。
\end{proof}

\begin{corollary}[闭环 Jacobi phase telescoping]
设 $\chi_1,\ldots,\chi_n\in\Gh$，令 $\chi_{n+1}=\chi_1$，并假设每个
$\chi_j$ 与 $\chi_j^{-1}\chi_{j+1}$ 及其乘积均非主。则
\[
\prod_{j=1}^n
\jmath(\chi_j,\chi_j^{-1}\chi_{j+1})
=\prod_{j=1}^n\rho(\chi_j^{-1}\chi_{j+1}).
\tag{7.5}
\]
特别地，所有依赖中间顶点 $\rho(\chi_j)$ 的 phase 精确相消。
\end{corollary}

\begin{proof}
由 (7.3)，第 $j$ 项为
$\rho(\chi_j)\rho(\chi_j^{-1}\chi_{j+1})
\overline{\rho(\chi_{j+1})}$。沿闭环相乘即可。
\end{proof}

\begin{remark}
若某个角色或乘积为主角色，$|\tau(1)|=1$ 而非 $\sqrt p$，(7.3) 的单位相位
版本失效。主角色、$t=1$、zero orbit 必须单独拆出；不能用 $p^{o(1)}$ 吞掉。
\end{remark}

\section{Gauss-square 的 multiplicative Fourier 对偶}

令
\[
\sigma(\chi)=\frac{\tau(\chi)^2}{p}
\quad(\chi\in\Gh),
\tag{8.1}
\]
这里包括主角色，且 $\tau(1)=-1$。

\begin{theorem}[Gauss-square $\leftrightarrow$ Kloosterman]
对每个 $t\in\G$，
\[
\boxed{
\sum_{\chi\in\Gh}\sigma(\chi)\chi(t)
=\frac Qp S(1,t^{-1};p).}
\tag{8.2}
\]
若 $\Kl_p(a)=p^{-1/2}S(1,a;p)$，则右侧为
\((Q/\sqrt p)\Kl_p(t^{-1})\)。
\end{theorem}

\begin{proof}
展开并交换求和：
\begin{align*}
\sum_\chi\tau(\chi)^2\chi(t)
&=\sum_{x,y\in\G}\ep(x+y)
\sum_\chi\chi(xyt)\\
&=Q\sum_{x\in\G}\ep\!\left(x+\frac1{xt}\right)
=Q S(1,t^{-1};p).
\end{align*}
除以 $p$ 即得结论。
\end{proof}

\begin{remark}
这是全角色族的 exact Fourier identity，不是对单个 $\chi$ 的点态估计。
其价值在于把一个高阶 Gauss phase family降回标准 rank-$2$ Kloosterman geometry；
但任何后续高矩估计仍需核对其 coefficient 与 support。
\end{remark}

\section{centered Parseval：主模应该精确删除}

\begin{theorem}[unitary centered Parseval]
对 $f:\F_p\to\C$ 及 (1.1) 的 Fourier convention，
\[
\boxed{
\sum_{h\ne0}|\widehat f(h)|^2
=\sum_{x\in\F_p}|f(x)|^2
-\frac1p\left|\sum_{x\in\F_p}f(x)\right|^2.}
\tag{9.1}
\]
\end{theorem}

\begin{proof}
Parseval 给 $\sum_h|\widehat f(h)|^2=\sum_x|f(x)|^2$，且
$\widehat f(0)=p^{-1/2}\sum_xf(x)$。减去 $h=0$ 项即得。
\end{proof}

\begin{corollary}[非单位化版本]
若 $\widehat f_{\rm raw}(h)=\sum_xf(x)\ep(-hx)$，则
\[
\boxed{
\sum_{h\ne0}|\widehat f_{\rm raw}(h)|^2
=p\sum_x|f(x)|^2-\left|\sum_xf(x)\right|^2.}
\tag{9.2}
\]
\end{corollary}

\begin{remark}
日志中多次出现的
$p\sum|Q|^2-|\sum Q|^2$ 正是 (9.2)，不是“额外得到一次 $p$ saving”。
它是删除 principal additive mode 的 exact identity；若 Fourier convention 改为
unitary，两个 $p$ 因子必须同步改变。
\end{remark}

\section{fibrewise dilation 与 Schatten transport}

令
\[
\mathscr K=\bigoplus_{R\in\G}\ell^2(\F_p).
\]

\begin{theorem}[charge-dependent additive dilation是酉的]
定义
\[
(\mathcal U f)(R,A)=f(R,A/R).
\tag{10.1}
\]
则 $\mathcal U$ 是 $\mathscr K$ 上的酉算子，逆为
\((\mathcal U^{-1}f)(R,A)=f(R,AR)\)。
\end{theorem}

\begin{proof}
固定 $R\in\G$ 时，$A\mapsto A/R$ 是 $\F_p$ 的置换，故保持该 fibre 的
$\ell^2$ 内积。对 $R$ 作正交直和即得。
\end{proof}

\begin{remark}
必须排除 $R=0$ 或为它另定义一个独立 fibre；在 $R=0$ 上写 $A/R$ 是未定义的。
此外，(10.1) 只证明该 dilation 本身酉，不证明它可以穿过任意 joint transform。
后者需要 exact intertwining identity。
\end{remark}

\begin{theorem}[Schatten ideal property]
设 $1\le s\le\infty$，$T$ 为有限维算子，$A,B$ 为有界算子。则
\[
\boxed{
\|ATB\|_{S_s}\le\|A\|_{op}\,\|T\|_{S_s}\,\|B\|_{op}.}
\tag{10.2}
\]
特别地，正交投影和酉算子不会增大 Schatten-$s$ 范数。
\end{theorem}

\begin{proof}
有限维 Schatten classes 是双边 operator ideals；可由 singular-value
不等式 $s_j(ATB)\le\|A\|_{op}\|B\|_{op}s_j(T)$ 后对 $j$ 取 $\ell^s$ 范数得到。
\end{proof}

\begin{corollary}
若某 exact factorization 为 $S=P_0U\Pi\mathcal AT$，其中 $P_0,\Pi$ 是
正交投影、$U$ 酉，且 $\mathcal A$ 为 (4.2)，则
\[
\|S\|_{S_s}\le\sqrt2\,\|T\|_{S_s}.
\tag{10.3}
\]
\end{corollary}

\begin{remark}
(10.3) 是条件结论：必须先逐项证明 exact factorization。仅凭“每个组成算子
看起来像 contraction”不能推断原算子等于它们的乘积。
\end{remark}

\section{为什么一般 entry shear 不保持 Schatten 范数}

\begin{proposition}[entry multiset 相同不保证 singular values 相同]
存在两个 $2\times2$ 矩阵，它们的 entries 只是彼此重排，但 Schatten-$4$ 范数不同。
\end{proposition}

\begin{proof}
取
\[
A=\begin{pmatrix}1&1\\0&0\end{pmatrix},\qquad
B=\begin{pmatrix}1&0\\0&1\end{pmatrix}.
\]
二者都有两个 $1$ 和两个 $0$。但 $A$ 的 singular values 为 $(\sqrt2,0)$，
$B$ 的 singular values 为 $(1,1)$。故
\[
\|A\|_{S_4}^4=4,
\qquad
\|B\|_{S_4}^4=2.
\]
\end{proof}

\begin{corollary}[合法坐标变换判据]
要无损运输 Schatten norm，足够的形式是 $T\mapsto UTV$，其中 $U,V$ 酉；
一般依赖 row 的 column shear、依赖 column 的 row shear或任意 entry bijection
不在此列，必须另证 intertwining identity。
\end{corollary}

\section{乘法能量与 Mellin Parseval}

对 $a:\G\to\C$ 定义乘法自相关
\[
C_a(t)=\sum_{x\in\G}a(tx)\overline{a(x)}.
\tag{12.1}
\]

\begin{theorem}[乘法 Wiener--Khinchin]
用 (1.2) 的 unitary Mellin convention，有
\[
\boxed{
(\mathcal M C_a)(\chi)=\sqrt Q\,|\mathcal Ma(\chi)|^2.}
\tag{12.2}
\]
从而
\[
\boxed{
\sum_{t\in\G}|C_a(t)|^2
=Q\sum_{\chi\in\Gh}|\mathcal Ma(\chi)|^4.}
\tag{12.3}
\]
\end{theorem}

\begin{proof}
展开左侧：
\begin{align*}
(\mathcal MC_a)(\chi)
&=Q^{-1/2}\sum_{t,x}a(tx)\overline{a(x)}\overline{\chi(t)}\\
&=Q^{-1/2}\sum_{u,x}a(u)\overline{a(x)}
\overline{\chi(u/x)}\\
&=\sqrt Q\,(\mathcal Ma)(\chi)\overline{(\mathcal Ma)(\chi)}.
\end{align*}
再用 Mellin Parseval 得 (12.3)。
\end{proof}

\begin{remark}
这是 exact $L^4$/energy dictionary。它可以把 coefficient fourth moment改写为
乘法碰撞，但不会自动控制带额外 projective/Kloosterman recoupler 的 Schatten cycle。
\end{remark}

\section{GCD-carried FSAC：共享素因子的严格修正版}

\begin{theorem}[primitive/GCD-carried factor-scale allocation]\label{thm:gcd-fsac}
设 $\mu(c)\mu(d)\ne0$，令
\[
g=(c,d),\qquad c=ga,\qquad d=gb.
\tag{13.1}
\]
则 $g,a,b$ 两两互素，且
\[
\mu(c)\mu(d)=\mu(a)\mu(b).
\tag{13.2}
\]
固定实数 $t_1,t_2$，置
\[
u=\frac{t_1+t_2}{2},\qquad v=\frac{t_1-t_2}{2}.
\]
若 $m=ab$，则
\[
\sum_{\substack{ab=m\\(a,b)=1\\(ab,g)=1}}
\mu(a)\mu(b)a^{-it_1}b^{-it_2}
=
\mathbf1_{(m,g)=1}\mu^2(m)m^{-iu}
\prod_{q\mid m}\bigl(-2\cos(v\log q)\bigr).
\tag{13.3}
\]
\end{theorem}

\begin{proof}
$c,d$ squarefree。共同素因子全部恰出现于 $g$，故 $g,a,b$ 两两互素。
于是
\[
\mu(c)\mu(d)=\mu(g)^2\mu(a)\mu(b)=\mu(a)\mu(b).
\]
又
\[
a^{-it_1}b^{-it_2}=(ab)^{-iu}(a/b)^{-iv}.
\]
因 $m$ squarefree，每个 $q\mid m$ 恰可分配到 $a$ 或 $b$。若分到 $a$，
局部贡献为 $-q^{-iu-iv}$；若分到 $b$，贡献为 $-q^{-iu+iv}$。
两者相加为 $-2q^{-iu}\cos(v\log q)$；对所有 $q\mid m$ 相乘即得。
\end{proof}

\begin{corollary}[reciprocal-zeta factorization]\label{cor:gcd-zeta}
定义
\[
A^{(g)}_{u,v}(m)=
\mathbf1_{(m,g)=1}\mu^2(m)m^{-iu}
\prod_{q\mid m}\bigl(-2\cos(v\log q)\bigr).
\tag{13.4}
\]
对 $\Re s>1$，
\[
\sum_{m\ge1}\frac{A^{(g)}_{u,v}(m)}{m^s}
=\frac{H_g(s;u,v)}
{\zeta(s+i(u-v))\zeta(s+i(u+v))},
\tag{13.5}
\]
其中除有限个 $q\mid g$ 的 Euler factors 外，$H_g$ 在 $\Re s>1/2$
绝对收敛；而这些有限 factors 的总大小在 $s$ 接近 $1$ 时至多为
$(g/\varphi(g))^{O(1)}$。
\end{corollary}

\begin{proof}
令
\[
x_q=q^{-s-i(u-v)},\qquad y_q=q^{-s-i(u+v)}.
\]
当 $q\nmid g$ 时，左侧局部因子为 $1-x_q-y_q$，故
\[
H_{g,q}=\frac{1-x_q-y_q}{(1-x_q)(1-y_q)}
=1+O(q^{-2\Re s}).
\]
当 $q\mid g$ 时，局部因子为 $1$，所以
$H_{g,q}=((1-x_q)(1-y_q))^{-1}$。前一类 factors 在 $\Re s>1/2$
绝对收敛；后一类有限，且由
$\prod_{q\mid g}(1-q^{-1})^{-O(1)}=(g/\varphi(g))^{O(1)}$ 控制。
\end{proof}

\begin{remark}
原日志中不带 $g$ 的 FSAC 公式只在 $(c,d)=1$ 时字面成立。定理
\ref{thm:gcd-fsac} 表明共享素因子不会摧毁 reciprocal-zeta deletion，
但必须作为独立 square/gcd carrier 保留，不能被 $\mu^2(m)$ 静默删除。
\end{remark}

\section{Origin-labelled tensorization 与区间自相关}

\begin{theorem}[typed packet reconstruction]\label{thm:typed-packet}
设
\[
a=\prod_{j=1}^Jm_{j,a},\quad b=\prod_{j=1}^Jm_{j,b},\quad
e=\prod_{j=1}^Jm_{j,e},\quad f=\prod_{j=1}^Jm_{j,f},
\tag{14.1}
\]
且不同 $j$ 保存 disjoint prime occurrences，$\mu(a)\mu(e)\ne0$。则
\[
\boxed{
\mu(a)\mu(e)\log b
=\sum_{j=1}^{J}
\left[\mu(m_{j,a})\mu(m_{j,e})\log m_{j,b}\right]
\prod_{i\ne j}\mu(m_{i,a})\mu(m_{i,e}).}
\tag{14.2}
\]
所以重新分桶后的 coefficient 不是“忘掉来源后的单变量函数”，而是秩至多
$J$ 的精确 typed tensor。
\end{theorem}

\begin{proof}
squarefree support 与 disjoint occurrences 给
$\mu(a)=\prod_j\mu(m_{j,a})$、$\mu(e)=\prod_j\mu(m_{j,e})$；
再展开 $\log b=\sum_j\log m_{j,b}$。
\end{proof}

\begin{corollary}[smooth monomial constraints 的 Fourier--Mellin 分离]
若 $\Psi\in C_c^\infty(\mathbb R^d)$ 作用于
$(\log m_1,\ldots,\log m_d)$，则
\[
\Psi(\log m_1,\ldots,\log m_d)
=\int_{\mathbb R^d}\widehat\Psi(\mathbf t)
\prod_{\nu=1}^dm_\nu^{it_\nu}\,d\mathbf t.
\tag{14.3}
\]
因此固定次数的 smooth support coupling 可精确 tensorize，总 cost 为
$\|\widehat\Psi\|_{L^1}$。
\end{corollary}

\begin{proof}
这是 $\mathbb R^d$ 上的 Fourier inversion。光滑紧支撑使
$\widehat\Psi$ 快速衰减，故绝对可积并可使用 Fubini。
\end{proof}

\begin{theorem}[interval autocorrelation 是 triangular weight]
设 $I\subset\mathbb Z$ 是长度 $H<p/2$ 的区间，
\[
w(x)=H^{-1/2}\mathbf1_I(x),\qquad
\Phi(s)=\sum_xw(x)\ep(sx).
\]
则
\[
|\Phi(s)|^2=\sum_vY(v)\ep(sv),\qquad
Y(v)=\left(1-\frac{|v|}{H}\right)_+
\tag{14.4}
\]
（把 $v$ 取为唯一满足 $|v|<p/2$ 的代表）。特别地，
\[
\supp Y\subset[-H,H],\qquad \|Y\|_\infty\le1,\qquad
\operatorname{TV}(Y)=2.
\tag{14.5}
\]
\end{theorem}

\begin{proof}
展开
\[
|\Phi(s)|^2
=\sum_{x,y}w(x)\overline{w(y)}\ep(s(x-y))
=\sum_v\left(\sum_yw(y+v)\overline{w(y)}\right)\ep(sv).
\]
两个长度 $H$ 区间相距 $v$ 时的交集大小为 $(H-|v|)_+$，故得到
(14.4)。其斜率在两侧分别为 $\pm H^{-1}$，总变差为 $2$。
\end{proof}

\begin{remark}
这说明 Burgess/short-window 步骤中常见的 triangular $Y$ 不是额外平坦性
假设，而是 normalized interval projection 的精确自相关；若换成任意权，
则必须重新证明其 support、$\ell^\infty$ 与 variation bounds。
\end{remark}

\section{Cross-prime character incidence：outer tensor leg 的正交}

\begin{theorem}[双尺度 cross-prime incidence operator]\label{thm:tool-crossprime}
设 $3\le H_2\le H_1$，$\mathcal P\subset[H_1,2H_1]$、
$\mathcal R\subset[H_2,2H_2]$ 为素数集合，并排除
$p=r$。在
\[
\bigoplus_{p\in\mathcal P}\ell^2(\widehat{\F_p^\times})
\quad\hbox{与}\quad
\bigoplus_{r\in\mathcal R}\ell^2(\widehat{\F_r^\times})
\]
之间定义
\[
\mathscr U_{(p,\chi),(r,\psi)}=\chi(r)^2\psi(p)^2.
\tag{15.1}
\]
则
\[
\boxed{\|\mathscr U\|_{2\to2}
\ll H_1H_2^{1/2}L^{1/2}},\qquad
L=1+\frac{\log(4H_1)}{\log H_2}.
\tag{15.2}
\]
对任意 Hilbert space $\mathcal H$，
$\|\mathscr U\otimes I_{\mathcal H}\|=\|\mathscr U\|$。特别地，
$H_1=H_2=H$ 时得到 $\|\mathscr U\|\ll H^{3/2}$。
\end{theorem}

\begin{proof}
令 $G=\mathscr U\mathscr U^*$。角色正交给
\[
G_{(p,\chi),(p',\chi')}
=\sum_{\substack{r\in\mathcal R,\ r\ne p,p'\\r\mid p^2-p'^2}}
(r-1)\chi(r)^2\overline{\chi'(r)^2}.
\tag{15.3}
\]
当 $p\ne p'$ 时，$r\mid(p-p')(p+p')$；乘积大小至多 $16H_1^2$，
故只有 $O(L)$ 个 $H_2$-scale prime $r$。每个 $r$ 产生秩一 block
$(r-1)v_{p,r}v_{p',r}^*$，而
$\|v_{p,r}\|_2=\sqrt{p-1}\ll H_1^{1/2}$，所以
\[
\|G_{p,p'}\|\ll LH_1H_2.
\tag{15.4}
\]
当 $p=p'$ 时，置 $V_{\chi,r}=\chi(r)^2$。再用角色正交，
\[
(V^*V)_{r,r'}=(p-1)\mathbf1_{r^2\equiv r'^2\pmod p}.
\tag{15.5}
\]
对固定 $r$，$r'\equiv\pm r\pmod p$ 在长度 $H_2\le H_1\le p$ 的区间中
各只有 $O(1)$ 个代表，故 Schur test 给 $\|V\|^2\ll H_1$。于是
\[
\|G_{p,p}\|
=\|V\diag(r-1)V^*\|\ll H_1H_2.
\tag{15.6}
\]
最后以至多 $O(H_1)$ 个 prime indices 作 block Schur：
\[
\|G\|\le\sup_p\sum_{p'}\|G_{p,p'}\|\ll LH_1^2H_2.
\]
取平方根即得。tensor identity 是 Hilbert 空间算子范数的标准性质。
\end{proof}

\begin{remark}
该定理只控制已经严格变成 (15.1) 的 outer character leg。CRT Gauss
factorization确实产生 $\chi(r)^2\psi(p)^2$，但真实算术 leaf 是否无损变成
$\mathscr U\otimes I_{\mathcal H}$ 仍需 basis、measure、prefactor 和 exceptional
characters 的逐项 compiler；不能只凭 (15.2) 宣布整个 P2 block 已证。对完整
dyadic blocks，朴素 aggregation scale 为
$H_1H_2(H_1H_2)^{o(1)}$，所以 (15.2) 显示 saving 由较小 prime scale
$H_2$ 决定，量级为 $H_2^{-1/2}$（允许对数因子）。
\end{remark}

\section{BP trace pushforward：sum--product 坐标与 divisor rigidity}

\begin{theorem}[trace polynomial 的双 sum--product 化]\label{thm:tool-bptrace}
对整数 $a,b,c,d$ 定义
\[
R(a,b,c,d)=abcd-(a+c)(b+d).
\tag{16.1}
\]
令
\[
x=ac,\quad y=a+c,\qquad u=bd,\quad v=b+d.
\]
则
\[
\boxed{R=xu-yv.}
\tag{16.2}
\]
映射 $(a,c)\mapsto(x,y)$ 的有序 fibre 至多二。固定 $a,c,R$，令
$A=ac,C=a+c$，则
\[
\boxed{(Ab-C)(Ad-C)=AR+C^2.}
\tag{16.3}
\]
\end{theorem}

\begin{proof}
(16.2) 由展开得到。给定 $(x,y)$ 时，$a,c$ 是
$T^2-yT+x$ 的有序根，故至多交换二重性。最后
\[
(Ab-C)(Ad-C)
=A^2bd-AC(b+d)+C^2
=A(Abd-C(b+d))+C^2=AR+C^2.
\]
\end{proof}

\begin{corollary}[严格可得的粗 pushforward energy]\label{cor:tool-crude}
设 $a,b,c,d\in[H,2H]$，并令
\[
a_R=\sum_{R(a,b,c,d)=R}
Y_1(a)Y_2(b)\overline{Y_3(c)Y_4(d)}.
\tag{16.4}
\]
则
\[
\boxed{
\sum_R|a_R|^2\ll_\varepsilon H^{2+\varepsilon}
\prod_{j=1}^4\|Y_j\|_2^2.}
\tag{16.5}
\]
\end{corollary}

\begin{proof}
固定 $a,c,R$ 后，(16.3) 把 $(b,d)$ 注入非零整数 $AR+C^2$ 的有序
因子分解；因 $A\asymp H^2,C\asymp H,b,d\asymp H$，两个左因子均非零。
故对固定 $(a,c,R)$ 有 $O_\varepsilon(H^\varepsilon)$ 个 $(b,d)$，
从而每个完整 $R$-fibre 至多 $H^{2+\varepsilon}$。对每个 fibre 用
Cauchy，再对 $R$ 求和，得到 (16.5)。
\end{proof}

\begin{remark}[严格适用范围]
(16.5) 对 bounded weights 给 $H^{6+\varepsilon}$，而最优 BPTPE 希望删除
前面的 $H^2$。外部 quadratic square-paucity 只数裸 $R$ 值，不能自动吸收
$a_R$ 的 pushforward multiplicity；weighted square-core BPTPE 因而仍是开放
analytic target。
\end{remark}

\section{Exact Parity--Gauge Decomposition}

令 $h(n)=(-1)^{\Omega(n)}$，Dirichlet convolution 记为 $*$，并置记逐点乘积。

\begin{theorem}[任意 divisor gauge 的精确 parity 分解]\label{thm:tool-epg}
对任意算术函数 $\lambda$，定义
\[
L_\lambda=\lambda*1,qquad
D_\lambda=\mu^2*(h\log)-(\lambda h)*h.                \tag{17.1}
\]
则对所有整数 $n\ge1$，
\[
\boxed{\Lambda(n)=L_\lambda(n)+h(n)D_\lambda(n).}     \tag{EPG}
\]
\end{theorem}

\begin{proof}
$h$ 完全乘法、$h^2=1$，且 $\mu h=\mu^2$。由 $\Lambda=\mu*\log$，
\[
h(n)\Lambda(n)
=\sum_{d\mid n}\mu(d)h(d)h(n/d)\log(n/d)
=(\mu^2*(h\log))(n).
\]
另一方面，
\[
h(n)L_\lambda(n)
=\sum_{d\mid n}\lambda(d)h(d)h(n/d)
=((\lambda h)*h)(n).
\]
相减并使用 $h^2=1$ 即得。
\end{proof}

\begin{corollary}[理想 prime-producing gauge]
令 $\lambda^\star(d)=-\mu(d)\log d$。则
\[
\lambda^\star=\mu*\Lambda,qquad
L_{\lambda^\star}=\Lambda,qquad D_{\lambda^\star}=0. \tag{17.2}
\]
\end{corollary}

\begin{proof}
$\sum_{d\mid n}\mu(d)\log d=-\Lambda(n)$，故
$1*\lambda^\star=\Lambda$。Möbius inversion给
$\lambda^\star=\mu*\Lambda$；代回 (EPG) 得 $D_{\lambda^\star}=0$。
\end{proof}

\begin{remark}
$h=e^{i\pi\Omega}$；在 squarefree support 上 $h=\mu$，且恒等式
$\mu^2h=\mu$ 对所有整数成立。因此 EPG 精确定位单一 parity spectral point，
但不估计 $D_\lambda$。理想 gauge 的困难也正是它包含完整 prime detector；
不能以“存在 $\lambda^\star$”冒充已经构造了可分析的 sieve weight。
\end{remark}

\section{finite ModelGlue 与 proof-compiler 守恒律}

\begin{theorem}[有限 exact ModelGlue]\label{thm:tool-modelglue}
令 $U,V,W$ 为复向量空间，$\mathfrak M:U\times V\to W$ 双线性，
$\widetilde{\mathfrak M}:U\otimes V\to W$ 为线性化。设 $T,S,\Sigma$ 有限，
\[
u=\sum_{t\in T}\kappa_tu_t,qquad
v=\sum_{s\in S}\ell_sv_s,qquad
\sum_{\sigma\in\Sigma}P_\sigma=I_{U\otimes V}.       \tag{18.1}
\]
置 $\mathfrak M_\sigma(a,b)=
\widetilde{\mathfrak M}(P_\sigma(a\otimes b))$。则
\[
\boxed{
\sum_{\sigma,t,s}\kappa_t\ell_s
\mathfrak M_\sigma(u_t,v_s)=\mathfrak M(u,v).}        \tag{MG}
\]
\end{theorem}

\begin{proof}
有限线性性把左端化为
\[
\widetilde{\mathfrak M}\!\left(
\Bigl(\sum_\sigma P_\sigma\Bigr)(u\otimes v)\right)
=\widetilde{\mathfrak M}(u\otimes v)=\mathfrak M(u,v).
\]
\end{proof}

\begin{corollary}[VDHGT/GFT synthesis 的合法版本]
若 $A=1+B$、$A'=1+B'$，则双线性性给
\[
\mathfrak M(A,A')=\mathfrak M(1,1)+\mathfrak M(B,1)
+\mathfrak M(1,B')+\mathfrak M(B,B').                 \tag{18.2}
\]
对绝对 Bochner 可积的连续 synthesis，先对简单函数/有限截断应用
定理~\ref{thm:tool-modelglue}，再用 dominated convergence 与 Fubini，可得同一结论。
\end{corollary}

\begin{proposition}[pure repartition cannot create a property]
令 $X,\Sigma$ 有限，$a_x\in\C$，并对每个 $x$ 有
$\sum_\sigma\rho_\sigma(x)=1$。若 $a_x\ne0$，则至少一项
$\rho_\sigma(x)a_x\ne0$。因此不改变原 term 的 exact partition 不能让一个
bad/nonrough term 从所有 descendants 中消失。
\end{proposition}

\begin{proof}
否则由 $a_x\ne0$ 得每个 $\rho_\sigma(x)=0$，与总和为 $1$ 矛盾。
\end{proof}

\begin{warning}[守恒不等于 cancellation]
取 $U=V=W=\C$、$\mathfrak M(u,v)=uv$、$P_0=I$、$u=v=1$，
ModelGlue 输出仍为 $1$。因此它只证明 principal model 在 exact splitting 下守恒，
不产生 $p^{-\delta}$ saving、roughness 或 G2/G3 estimate。古早 raw BAL-2P
被 BAL-SPARSE 覆盖的结论正违反上述 no-go；加入 least-prime/Buchstab identity
才是改变 terms 的合法新步骤。
\end{warning}

\section{factor-count 与 factor-scale 的完整有限谱}

固定 $X>1$。对 $r\ge1$ 定义
\[
P_r^{(X)}(n)=\sum_{p^a\Vert n}a
\left(\frac{\log p}{\log X}\right)^r,qquad
P_0(n)=\Omega(n).                                     \tag{19.1}
\]

\begin{theorem}[完全加性与 unitary factor-spectral phase]\label{thm:tool-gfst}
对所有正整数 $m,n$ 及 $r\ge0$，
\[
P_r^{(X)}(mn)=P_r^{(X)}(m)+P_r^{(X)}(n).              \tag{19.2}
\]
因此对 $R\ge1$、$\vartheta,t_1,\ldots,t_R\in\mathbb R$，
\[
\chi_{\vartheta,\mathbf t}^{(X)}(n)
=\exp\!\left(i\vartheta\Omega(n)+i\sum_{r=1}^Rt_rP_r^{(X)}(n)\right)
\tag{19.3}
\]
是 modulus-one completely multiplicative function。特别地，
$\chi_{\pi,\mathbf0}=h$，且 $\mu^2\chi_{\pi,\mathbf0}=\mu$。
\end{theorem}

\begin{proof}
对每个素数 $p$，$v_p(mn)=v_p(m)+v_p(n)$；代入 (19.1) 后逐素数求和即得
(19.2)，包括 $r=0$。指数把加法变乘法，证明 (19.3)。最后
$e^{i\pi\Omega}=(-1)^\Omega$；squarefree 时它等于 $\mu$，非 squarefree
时 $\mu^2=0$，故最后一式对所有整数成立。
\end{proof}

\begin{theorem}[rough support 上的 Newton 重构]\label{thm:tool-newton}
设 $0<\delta\le1$、$n\le X$，且每个 $p\mid n$ 满足 $p\ge X^\delta$。
令 $k=\Omega(n)$、$K=\lfloor1/\delta\rfloor$。则 $k\le K$，并且有限数据
\[
\boxed{\bigl(k,P_1^{(X)}(n),\ldots,P_K^{(X)}(n)\bigr)} \tag{19.4}
\]
唯一确定计重数的 log-scale multiset
$\{\log p/\log X:p\mid n\}$。
\end{theorem}

\begin{proof}
把计重数的素因子写成 $p_1,\ldots,p_k$，并置
$u_i=\log p_i/\log X\ge\delta$。因 $n\le X$，
\[
1\ge\frac{\log n}{\log X}=\sum_{i=1}^ku_i\ge k\delta,
\]
故 $k\le\lfloor1/\delta\rfloor=K$。而
$P_r=\sum_{i=1}^ku_i^r$ 是 $u_i$ 的 power sums。令 $e_0=1$，Newton identities
递归给
\[
re_r=\sum_{j=1}^r(-1)^{j-1}e_{r-j}P_j,qquad1\le r\le k. \tag{19.5}
\]
所以 (19.4) 决定 $e_1,\ldots,e_k$，进而决定首一多项式
\[
T^k-e_1T^{k-1}+e_2T^{k-2}-\cdots+(-1)^ke_k,
\]
其根的 multiset 正是 $\{u_1,\ldots,u_k\}$。
\end{proof}

\begin{remark}
这是 exact finite reconstruction，不是稳定性定理；近似 moments 的数值反演可能
病态。尺度谱即使完全重构 factor geometry，也不会自动控制 parity spectral point
$(\vartheta,\mathbf t)=(\pi,\mathbf0)$。所以“generic $\mathbf t$ modes 小”不能替代
Liouville/Möbius cancellation。
\end{remark}

\section{finite alternating jet：endpoint 的精确 Newton 公式}

\begin{lemma}[截断二项式恒等式]\label{lem:tool-binomial}
对整数 $k,M\ge0$，
\[
\sum_{j=0}^M(-1)^j\binom kj
=\begin{cases}
1,&k=0,\\
0,&1\le k\le M,\\
(-1)^M\binom{k-1}{M},&k>M.
\end{cases}                                           \tag{20.1}
\]
\end{lemma}

\begin{proof}
$k=0$ 直接成立。对 $k\ge1$，Pascal 恒等式给
\[
\sum_{j=0}^M(-1)^j\binom kj
=\sum_{j=0}^M(-1)^j\binom{k-1}j
-\sum_{j=0}^{M-1}(-1)^j\binom{k-1}j
=(-1)^M\binom{k-1}M.
\]
当 $1\le k\le M$ 时最后的二项式系数为零。
\end{proof}

\begin{theorem}[有限 endpoint extraction]\label{thm:tool-jet}
令 $A$ 为有限集，$K:A\to\mathbb N$，$w:A\to\C$，并定义
\[
A_0=\sum_{K(a)=0}w(a),qquad
J_j=\sum_{a\in A}w(a)\binom{K(a)}j,qquad
T_M=\sum_{K(a)>M}w(a)\binom{K(a)-1}M.                 \tag{20.2}
\]
则对每个 $M\ge0$，
\[
\boxed{A_0=\sum_{j=0}^M(-1)^jJ_j+(-1)^{M+1}T_M.}     \tag{JET}
\]
\end{theorem}

\begin{proof}
交换两个有限和，并对每个 $a$ 应用引理~\ref{lem:tool-binomial}。当 $K(a)=0$
时贡献为 $w(a)$；$1\le K(a)\le M$ 时为零；$K(a)>M$ 时由尾项精确抵消。
\end{proof}

\begin{corollary}[Goldbach endpoint 的纯组合版本]
固定偶数 $N$，在任意满足 $2\le n\le N-2$ 的有限 support 上置
\[
K_N(n)=\Omega(n)+\Omega(N-n)-2.                       \tag{20.3}
\]
则 $K_N(n)\ge0$，且
\[
K_N(n)=0\iff n\ \text{与}\ N-n\ \text{都是素数}.   \tag{20.4}
\]
因此定理~\ref{thm:tool-jet} 对 $K_N$ 给出 weighted Goldbach count 的精确
finite-jet 加 tail 表示。
\end{corollary}

\begin{proof}
每个大于 $1$ 的整数满足 $\Omega\ge1$，且 $\Omega(m)=1$ 当且仅当 $m$ 为素数。
于是 (20.3)--(20.4) 成立，再调用 (JET)。
\end{proof}

\begin{warning}[ASJT 的真实逻辑强度]
(JET) 是有限组合恒等式；它没有估计 signed jet 或 tail。当 tail 已经小，精确到
$o(N/\log^2N)$ 的 signed-jet asymptotic 就在逼近 $K_N=0$ 极端下尾，不能因其
写在 principal point $z=1$ 附近就视为容易。古早 ASJT 路线仍缺这一 analytic
核心及独立的 tail/local-model interfaces。
\end{warning}

\section{RSC/Bellman：无限完成不能制造严格正性}

\begin{theorem}[strict positivity has a finite witness]\label{thm:tool-finite-witness}
设 $I$ 非空，$v_i\in\mathbb R$，且 $V=\sup_{i\in I}v_i$ 存在。则
\[
\boxed{V>0\iff\exists i\in I:\ v_i>0.}               \tag{21.1}
\]
\end{theorem}

\begin{proof}
反向显然。若不存在正的 $v_i$，则 $0$ 是全体 $v_i$ 的上界，故最小上界
$V\le0$，与 $V>0$ 矛盾。
\end{proof}

\begin{theorem}[monotone Bellman supersolution comparison]\label{thm:tool-bellman}
令 $Z$ 为集合，$\mathcal F:(Z\to\mathbb R)\to(Z\to\mathbb R)$ 单调。
设 $V_0\le\Phi$、$\mathcal F\Phi\le\Phi$，并递归定义
$V_{r+1}=\mathcal FV_r$。则对每个 $r\ge0$，$V_r\le\Phi$。若逐点 supremum
$V_\omega=\sup_rV_r$ 存在，则 $V_\omega\le\Phi$。
\end{theorem}

\begin{proof}
对 $r$ 归纳。初始步是假设；若 $V_r\le\Phi$，单调性给
$V_{r+1}=\mathcal FV_r\le\mathcal F\Phi\le\Phi$。取逐点 supremum得到最后一句。
\end{proof}

\begin{remark}
定理~\ref{thm:tool-finite-witness} 说明若 RSC completion 定义为所有有限证书值的
supremum，则 endpoint 严格正性一定已经由某个有限证书实现；无限深极限可以令
临界参数趋近 $1$，却不能在端点凭空产生正值。定理~\ref{thm:tool-bellman} 给
dual barrier comparison。要声称 strong duality 或 $V_\omega$ 是 fixed point，
还必须另证 order-continuity、feasibility 与相应 action-domain hypotheses。
\end{remark}

\begin{warning}[naive parity invariance 对 full RSC 为假]
若 full RSC grammar 允许 terminal/weight 直接查询
$\1_{\Omega(n)=1}$，则它显然能区分 $\Omega(n)=1$ 与 $\Omega(n)=2$ 的状态；
所以不能宣称所有结构动作天然 parity-blind。严格可用的 no-go 只能是条件性的
proof-transport 定理：必须逐项假设每个 structural identity 与每个 terminal
oracle 都能运输到 parity-adversarial model，再对有限 derivation tree 作结构归纳。
“语法能表达 prime”与“analytic oracle 能证明 prime branch 有正质量”是两件事。
\end{warning}

\section{外部定理：只能按原适用域调用}

本节不把外部论文重新证明为本文定理，只记录经原文核对的调用边界。

\begin{center}
\begin{longtable}{p{0.24\textwidth}p{0.32\textwidth}p{0.34\textwidth}}
\toprule
外部结果 & 原文可提供 & 不能自动推出\\
\midrule
\endfirsthead
\toprule
外部结果 & 原文可提供 & 不能自动推出\\
\midrule
\endhead
de la Vall\'ee--Poussin zero-free region
& $\zeta(s)$ 在 $\Re s\ge1-c/\log(|\Im s|+3)$ 的经典无零区域及
$1/\zeta$ 的标准多对数界
& RH；越过 zeta 零点的无条件大范围 contour shift\\
Montgomery--Vaughan singular tail
& 点态界 $|\widetilde{\mathfrak S}_Y(k)|
\ll d(k)(\log\log3Y)^2/Y$
& 在 $Y=\log^B N$ 时 uniform 的任意 log-power saving\\
Kable (2002)
& Legendre/Soto--Andrade 作为 Steinberg tensor-square 的 parabolic CG
coefficients；在 $\ell^2(\F_q,m)$ 中完备正交
& 我们 counting measure 上的任意 coefficient operator diagonal；DCCS；
未处理端点权的 unitary splice\\
Pascadi (2026), Prop. 5.1
& 有限群、正规子群、不可约表示上的 even Schatten amplification
& 任意 reducible block；忽略 character denominator；自动 outer-modulus saving\\
Pascadi (2026), Thm. 7.1
& 若 $c=dd'e$、$d'\mid d$、$(d,e)=1$，则以
$G=dMN(M^2+N^2)c^{-3}+f(M^2+N^2)c^{-2}+f/d^2$
给出 $G^{1/6}$ saving；$M,N$ 的对称表述见 Blomer--Pascadi Thm.~5.4；
balanced semiprime 可得约 $c^{-1/12}$
& prime modulus；未验证 coefficient/support/coprimality 的 G3 cell\\
Blomer--Pascadi (2026), Thm. 1.1
& 任意模数的 bilinear Kloosterman form；平方根临界区 saving $c^{-1/32}$
& 变化模数的 outer average；GFT weighted pushforward energy\\
de la Bret\`eche--Wang--Xu (2026)
& 固定可分离多项式值的未加权 square-paucity；二次情形最强
& 任意权 $a_R$ 的 BPTPE；$R=xu-yv$ pushforward 后的 weighted theorem\\
\bottomrule
\end{longtable}
\end{center}

\section{Lean-friendly theorem ledger}

\begin{center}
\begin{tabular}{p{0.45\textwidth}p{0.16\textwidth}p{0.27\textwidth}}
\toprule
结果 & 状态 & 主要类型/guard\\
\midrule
smooth HWAC/residue & 已完整证明 & smooth partition；$\ell^1$ Fourier\\
sharp/global-holomorphic HWAC & 已否定 & sharp Wiener norm 发散；无全纯延拓\\
Möbius zero-orbit resummation & 已完整证明 & $M,X,Y\asymp N$；zeta zero-free input\\
High-Conductor ZOT & 已完整证明 & smooth arcs；MV pointwise tail\\
有限 DFT projector & 已完整证明 & $0\le r,j<K$\\
Mellin inverse & 已完整证明 & $\G$，counting measure\\
reciprocal fibre & 已完整证明 & $z,z',y,y'\ne0$，$z\ne z'$\\
projective charge map & 已完整证明 & $q\ne-z$\\
joint Gram $I+S-D$ & 已完整证明 & exact $(R,P)$ 或 character $(R,\eta)$ 坐标\\
Vi\`ete diagonal $x+z=xz$ & 已完整证明 & $r\ne0,1$\\
Cayley--Kloosterman identity & 已完整证明 & 所有 $A,C\in\F_p$\\
Jacobi coboundary & 已完整证明 & $\alpha,\beta,\alpha\beta$ 非主\\
Gauss-square Fourier identity & 已完整证明 & 全部乘法角色\\
centered Parseval & 已完整证明 & 指定 Fourier normalization\\
fibrewise dilation unitary & 已完整证明 & $R\in\G$\\
Schatten ideal transport & 已完整证明 & 有限维\\
一般 shear no-go & 已完整证明 & 显式 $2\times2$ 反例\\
乘法 autocorrelation $L^4$ identity & 已完整证明 & 指定 Mellin normalization\\
GCD-carried FSAC & 已完整证明 & squarefree Möbius legs；共享 gcd 显式保留\\
typed packet reconstruction & 已完整证明 & origin occurrences 不可遗忘\\
interval autocorrelation & 已完整证明 & normalized interval，$H<p/2$\\
cross-prime incidence & 已完整证明 & 双 dyadic prime scales、全部角色、排除 $p=r$\\
BP trace/divisor rigidity & 已完整证明 & 整数短区间；只推出粗 $H^6$ energy\\
Exact Parity--Gauge Decomposition & 已完整证明 & 所有 $n\ge1$；任意算术 gauge\\
finite ModelGlue / repartition no-go & 已完整证明 & 有限 tensor partition；积分版需绝对可积\\
factor-scale Newton reconstruction & 已完整证明 & $n\le X$、$X^\delta$-rough、精确 moments\\
finite alternating endpoint jet & 已完整证明 & 任意有限 weighted set；不含 analytic tail bound\\
RSC finite witness / Bellman comparison & 已完整证明 & 实数 supremum；单调 operator\\
\bottomrule
\end{tabular}
\end{center}

\paragraph{形式化说明。}
上表中的有限域、有限群条目可在有限类型层形式化；Schatten 条目需要
singular-value theory，HWAC 需要 Fourier/Wiener 与复系数抽取，ZOT 还需要
Perron、Euler product、经典 zeta 零点自由区及 Montgomery--Vaughan tail 的
形式化封装。本文没有提供可编译 Lean 文件，故“Lean-friendly”不等于
“已由 Lean kernel 验证”。

\section{这些性质对下一步决策的含义}

\begin{enumerate}
\item HWAC 的真正复分析优势是把 smooth arc cutoff 变成绝对可和的 shifted
coefficient residues；不能把这一点夸张成 cutoff 本身的全纯延拓。
\item ZOT 表明 zero orbit 应延迟取绝对值：低 conductor 由 minor plateau
消去，高 conductor 与 major truncation 重组为完整奇异级数。
\item GFT 最稳定的优势不是“每个 mode 更小”，而是同时保留 additive sum 与
multiplicative product 后，root multiplicity 从表面上的 $p$ 级降到 swap 二重性。
\item Jacobi phase在非主闭环中是 coboundary，可被 gauge/telescoping；
真正困难通常在 coefficient placement 和 exceptional carriers，而非 phase 本身。
\item projective rational kernel可经 Cayley 变换精确回到经典 Kloosterman；
但 fixed-entry Weil 与 Schatten assembly必须分开证明。
\item centered subtraction 是正交投影，不是启发式误差删除；所有 principal/zero
mode 都应在 Fourier convention 下精确记账。
\item DCCS 的 fibrewise dilation本身已经是酉的；剩余问题只可能在它与真实
joint transform/short tensor的相对位置，而不是 dilation 的范数。
\item shared prime factors 应先进入 gcd carrier，再对互素残余使用 FSAC；
否则 $\mu^2(m)$ 会错误删除原和中的 $q^2$ 状态。
\item cross-prime incidence 确实给 canonical outer character leg 一个由较小
prime scale 决定的 $H_{\min}^{-1/2}$ operator saving（允许对数因子），但不会
替代 leaf-to-spectral compiler。
\item BP trace 的优美坐标是 $R=xu-yv$，且固定一对变量后有精确 divisor
factorization；当前它严格给出粗能量，weighted square-core saving 仍需新定理。
\item EPG 把任意 sieve gauge 的差精确钉在 $\vartheta=\pi$ carrier；它的价值是
定位 parity 信息，不是自动消去该信息。
\item ModelGlue 保证 exact compiler 不改变 principal term，而 repartition no-go
保证 bookkeeping 不能创造 roughness；新 analytic input 与新 identity 必须显式标记。
\item $P_0,\ldots,P_K$ 在 rough support 上完整重构 factor-scale multiset；因此继续
增加尺度 moments 只增加几何分辨率，不能替代 $(\pi,\mathbf0)$ parity control。
\item finite jet 精确抽取 endpoint，但 signed jet 的所需精度可能承载全部极端下尾
困难；它是替代表述，不是已经获得的 Goldbach saving。
\item RSC 的严格正值必有 finite witness；研究无限 completion 时应把“参数趋近端点”
与“端点出现正证书”分开，并用 supersolution 做下界审计。
\end{enumerate}

\begin{thebibliography}{9}
\bibitem{Titchmarsh}
E.~C. Titchmarsh, revised by D.~R. Heath-Brown,
\emph{The Theory of the Riemann Zeta-function}, 2nd ed., Oxford, 1986.

\bibitem{MVBook}
H.~L. Montgomery and R.~C. Vaughan,
\emph{Multiplicative Number Theory I: Classical Theory}, Cambridge, 2007.

\bibitem{GHN2014}
D.~A. Goldston, J. Ziegler Hunts and T. Ngotiaoco,
\emph{The Tail of the Singular Series for the Prime Pair and Goldbach Problems},
arXiv:1409.2151 (2014), especially equations (1.4), (1.8), (1.9).
\url{https://arxiv.org/abs/1409.2151}

\bibitem{Kable2002}
A.~C. Kable,
\emph{Legendre sums, Soto--Andrade sums and Kloosterman sums},
Pacific J. Math. 206 (2002), 139--157.
\url{https://msp.org/pjm/2002/206-1/pjm-v206-n1-p09-s.pdf}

\bibitem{Pascadi2026}
A. Pascadi,
\emph{Non-abelian amplification and bilinear forms with Kloosterman sums},
arXiv:2511.08445v2 (2026).
\url{https://arxiv.org/abs/2511.08445}

\bibitem{BP2026}
V. Blomer and A. Pascadi,
\emph{Bilinear forms with Kloosterman sums via quadratic characters},
arXiv:2607.24311v1 (2026).
\url{https://arxiv.org/abs/2607.24311}

\bibitem{dBWX2026}
R. de la Bret\`eche, V.~Y. Wang and M.~W. Xu,
\emph{Random Multiplicative Functions and Making Squares from Polynomial Values},
arXiv:2607.06398v1 (2026).
\url{https://arxiv.org/abs/2607.06398}
\end{thebibliography}

\end{document}
